\section{Methodology}
\label{sec:methodology}

This section describes how the study was conducted. The formal definitions of the multi-layer dependency graphs are presented separately in Appendix~\ref{sec:graph-definitions}. \label{sec:workflow} This paper is a human-AI collaborative open-source research project. We invite readers to regard it as a continuously evolving effort rather than a traditional, self-contained academic publication. All analyses were conducted in a command-line environment; the complete operation history, commit log, and revision trail are publicly available in the GitHub repository for inspection, reproduction, and critique. %The core research insight originates from the intrinsic connections within the language itself: \module{Mathlib}, as a formalized mathematical corpus, makes its internal dependency relations into a structural manifestation of the mathematical knowledge network. 
Our methodological starting point is to follow the internal logic of \module{Mathlib} and to deploy a toolchain for systematic extraction and analysis. Within this framework, humans drive the research questions, interpret results, and decide what to include or discard, while AI orchestrates and executes open-source tools (managing dependency installation, running extraction scripts, generating analysis code, and coordinating data flow between tools).

Raw data extraction relies on three complementary open-source tools: \texttt{importGraph}~\cite{importgraph} (module-level import edges), \texttt{jixia}~\cite{jixia2024} (declaration-level dependency edges and tactic usage profiles), and \texttt{lean-training-data}~\cite{lean_training_data} (declaration metadata). The resulting dataset comprises $317{,}655$ nodes, $8.4$M edges across $7{,}563$ modules, and is publicly available on HuggingFace. All analyses are implemented in Python using \texttt{NetworkX}~\cite{hagberg2008} (graph construction and centrality), \texttt{python-louvain}~\cite{aynaud2020} (community detection), \texttt{powerlaw}~\cite{alstott2014} (heavy-tailed fitting), \texttt{scikit-learn}~\cite{pedregosa2011} (NMI/ARI evaluation), \texttt{Pandas}~\cite{mckinney2010}, \texttt{NumPy}~\cite{harris2020}, and \texttt{Matplotlib}~\cite{hunter2007}. Analysis scripts follow a test-driven development cycle: write the test first, then implement, then verify incrementally.
 